% paper-skeleton.tex -- nfact-primorial-square paper lane (2026-10-06;
% disclosure cleanup 2026-10-07). This file is the SOURCE of paper.tex: the
% derived paper.tex is produced from it by audit/build.py, which substitutes
% the B:<name> markers with byte-exact contiguous excerpts of the shipped
% proof files and the N:<key> markers with figures measured at build time.
% The fact-closure gate scans this skeleton (not the derived file), so every
% fact-type figure here is a marker or a spelled-out word and no hand-copied
% number can go stale. Line ranges of the excerpts are recorded in
% audit/listing-verify.txt, not restated here.
% Since 2026-10-07 the assembly sources are the shipped comment-cleaned
% copies under zenodo/proofs/ (their code regions are byte-identical to the
% frozen development of record; correspondence record:
% zenodo/audit/comment-cleanup-diff.md).
% Companion claim of record: Zenodo concept DOI 10.5281/zenodo.23177003
% (published 2026-10-06).
\documentclass[11pt]{article}

\usepackage[T1]{fontenc}
\usepackage[utf8]{inputenc}
\usepackage{amsmath,amssymb,amsthm}
\usepackage{graphicx}
\usepackage{hyperref}
\usepackage{xurl}
\setlength{\emergencystretch}{3em}
% lstlean.tex — Lean style for the listings package.
% Lineage: Jeremy Avigad (2015), adapted from Assia Mahboubi's SSR style;
% inspired by Olivier Verdier's unixode.sty for the Unicode literate table.
% No CTAN package or upstream repo exists; papers carry this file in their
% source package (cap set arXiv:1907.01449, sphere eversion arXiv:2210.07746).
% Adaptation note: sorry/admit/sorryAx are highlighted in a dedicated error
% tier so that incomplete proofs stand out visually in listings.
\usepackage{listings}
\usepackage{xcolor}

\definecolor{keywordcolor}{rgb}{0.7, 0.1, 0.1}   % red keywords
\definecolor{tacticcolor}{rgb}{0.0, 0.1, 0.6}    % blue tactics
\definecolor{commentcolor}{rgb}{0.4, 0.4, 0.4}   % grey comments
\definecolor{symbolcolor}{rgb}{0.0, 0.1, 0.6}    % blue symbols
\definecolor{sortcolor}{rgb}{0.1, 0.5, 0.1}      % green sorts
\definecolor{errorcolor}{rgb}{0.6, 0, 0}         % dark red: sorry/admit
\definecolor{stringcolor}{rgb}{0.5, 0.1, 0.5}    % purple strings

\lstdefinelanguage{lean}{
  % Anything between $ becomes LaTeX math mode
  mathescape=false,
  % Comments may or not include Latex commands
  texcl=false,
  % keywords, list taken from lean-syntax.txt
  morekeywords=[1]{import, prelude, protected, private, noncomputable,
    definition, meta, renaming, hiding, parameter, parameters, begin,
    constant, constants, lemma, variable, variables, theory, print,
    theorem, example, open, section, namespace, universe, universes, end,
    modifier, modifiers, using, export, exposing, axiom, inductive,
    coinductive, with, structure, class, instance, attribute, local, where,
    by, have, show, let, in, if, then, else, match, do, for, fun, assume,
    from, take, calc, include, omit, infix, infixl, infixr, notation,
    macro, syntax, elab, set_option, initialize, deriving, opaque, abbrev,
    mutual, partial, scoped, termination_by, decreasing_by},
  % Sorts
  morekeywords=[2]{Sort, Type, Prop},
  % tactics, list taken from lean-syntax.txt (tactic keywords)
  morekeywords=[3]{intro, intros, apply, exact, refine, rfl, simp, simp_all,
    simpa, simp_rw, rw, rewrite, erewrite, ring, ring_nf, omega, decide,
    norm_num, norm_cast, induction, cases, rcases, obtain, constructor,
    ext, funext, conv, congr, field_simp, linarith, nlinarith, positivity,
    tauto, trivial, aesop, use, by_cases, by_contra, push_neg, contrapose,
    symm, trans, assumption, contradiction, exfalso, specialize,
    generalize, revert, clear, rename_i, subst, unfold, delta, change,
    convert, split_ifs, interval_cases, fin_cases, zify, gcongr,
    nth_rewrite, suffices, generalizing, at, only, native_decide},
  % warnings - constitution tier: must never survive into a final proof
  morekeywords=[4]{sorry, admit, sorryAx},
  literate=
    {α}{{\ensuremath{\alpha}}}1
    {β}{{\ensuremath{\beta}}}1
    {γ}{{\ensuremath{\gamma}}}1
    {δ}{{\ensuremath{\delta}}}1
    {ε}{{\ensuremath{\varepsilon}}}1
    {ζ}{{\ensuremath{\zeta}}}1
    {η}{{\ensuremath{\eta}}}1
    {θ}{{\ensuremath{\theta}}}1
    {ι}{{\ensuremath{\iota}}}1
    {κ}{{\ensuremath{\kappa}}}1
    {λ}{{\ensuremath{\lambda}}}1
    {μ}{{\ensuremath{\mu}}}1
    {ν}{{\ensuremath{\nu}}}1
    {ξ}{{\ensuremath{\xi}}}1
    {π}{{\ensuremath{\pi}}}1
    {ρ}{{\ensuremath{\rho}}}1
    {σ}{{\ensuremath{\sigma}}}1
    {τ}{{\ensuremath{\tau}}}1
    {υ}{{\ensuremath{\upsilon}}}1
    {φ}{{\ensuremath{\varphi}}}1
    {χ}{{\ensuremath{\chi}}}1
    {ψ}{{\ensuremath{\psi}}}1
    {ω}{{\ensuremath{\omega}}}1
    {Γ}{{\ensuremath{\Gamma}}}1
    {Δ}{{\ensuremath{\Delta}}}1
    {Θ}{{\ensuremath{\Theta}}}1
    {Λ}{{\ensuremath{\Lambda}}}1
    {Ξ}{{\ensuremath{\Xi}}}1
    {Π}{{\ensuremath{\Pi}}}1
    {Σ}{{\ensuremath{\Sigma}}}1
    {Φ}{{\ensuremath{\Phi}}}1
    {Ψ}{{\ensuremath{\Psi}}}1
    {Ω}{{\ensuremath{\Omega}}}1
    {ℵ}{{\ensuremath{\aleph}}}1
    {∀}{{\ensuremath{\forall}}}1
    {∃!}{{\ensuremath{\exists}!}}1
    {∃}{{\ensuremath{\exists}}}1
    {∈}{{\ensuremath{\in}}}1
    {∉}{{\ensuremath{\notin}}}1
    {∘}{{\ensuremath{\circ}}}1
    {∩}{{\ensuremath{\cap}}}1
    {∪}{{\ensuremath{\cup}}}1
    {⊆}{{\ensuremath{\subseteq}}}1
    {⊂}{{\ensuremath{\subset}}}1
    {⊇}{{\ensuremath{\supseteq}}}1
    {⊃}{{\ensuremath{\supset}}}1
    {⊄}{{\ensuremath{\nsubseteq}}}1
    {∅}{{\ensuremath{\emptyset}}}1
    {∧}{{\ensuremath{\wedge}}}1
    {∨}{{\ensuremath{\vee}}}1
    {¬}{{\ensuremath{\neg}}}1
    {⊤}{{\ensuremath{\top}}}1
    {⊥}{{\ensuremath{\bot}}}1
    {↔}{{\ensuremath{\leftrightarrow}}}1
    {→}{{\ensuremath{\rightarrow}}}1
    {←}{{\ensuremath{\leftarrow}}}1
    {↑}{{\ensuremath{\uparrow}}}1
    {⇒}{{\ensuremath{\Rightarrow}}}1
    {⇐}{{\ensuremath{\Leftarrow}}}1
    {⟹}{{\ensuremath{\Longrightarrow}}}1
    {↦}{{\ensuremath{\mapsto}}}1
    {≤}{{\ensuremath{\leq}}}1
    {≥}{{\ensuremath{\geq}}}1
    {≠}{{\ensuremath{\neq}}}1
    {≈}{{\ensuremath{\approx}}}1
    {≡}{{\ensuremath{\equiv}}}1
    {≃}{{\ensuremath{\simeq}}}1
    {≅}{{\ensuremath{\cong}}}1
    {×}{{\ensuremath{\times}}}1
    {·}{{\ensuremath{\cdot}}}1
    {±}{{\ensuremath{\pm}}}1
    {⊕}{{\ensuremath{\oplus}}}1
    {⊗}{{\ensuremath{\otimes}}}1
    {⋆}{{\ensuremath{\star}}}1
    {▸}{{\ensuremath{\blacktriangleright}}}1
    {⟨}{{\ensuremath{\langle}}}1
    {⟩}{{\ensuremath{\rangle}}}1
    {⦃}{{\ensuremath{\{\!\mid}}}1
    {⦄}{{\ensuremath{\mid\!\}}}}1
    {⟦}{{\ensuremath{[\![}}}1
    {⟧}{{\ensuremath{]\!]}}}1
    {⌊}{{\ensuremath{\lfloor}}}1
    {⌋}{{\ensuremath{\rfloor}}}1
    {⌈}{{\ensuremath{\lceil}}}1
    {⌉}{{\ensuremath{\rceil}}}1
    {√}{{\ensuremath{\surd}}}1
    {∣}{{\ensuremath{\mid}}}1
    {∤}{{\ensuremath{\nmid}}}1
    {∎}{{\ensuremath{\blacksquare}}}1
    {⋯}{{\ensuremath{\cdots}}}1
    {…}{{\ensuremath{\ldots}}}1
    {∑}{{\ensuremath{\sum}}}1
    {∏}{{\ensuremath{\prod}}}1
    {⁻¹}{{\ensuremath{^{-1}}}}1
    {¹}{{\ensuremath{^1}}}1
    {²}{{\ensuremath{^2}}}1
    {³}{{\ensuremath{^3}}}1
    {ⁿ}{{\ensuremath{^n}}}1
    {ᵐ}{{\ensuremath{^m}}}1
    {₀}{{\ensuremath{_0}}}1
    {₁}{{\ensuremath{_1}}}1
    {₂}{{\ensuremath{_2}}}1
    {₃}{{\ensuremath{_3}}}1
    {₄}{{\ensuremath{_4}}}1
    {₅}{{\ensuremath{_5}}}1
    {ₙ}{{\ensuremath{_n}}}1
    {ₘ}{{\ensuremath{_m}}}1
    {ₖ}{{\ensuremath{_k}}}1
    {ᵢ}{{\ensuremath{_i}}}1
    {ⱼ}{{\ensuremath{_j}}}1
    {ℤ}{{\ensuremath{\mathbb{Z}}}}1
    {ℕ}{{\ensuremath{\mathbb{N}}}}1
    {ℚ}{{\ensuremath{\mathbb{Q}}}}1
    {ℝ}{{\ensuremath{\mathbb{R}}}}1
    {ℂ}{{\ensuremath{\mathbb{C}}}}1
    {𝔽}{{\ensuremath{\mathbb{F}}}}1
    {𝒫}{{\ensuremath{\mathcal{P}}}}1
    {∗}{{\ensuremath{\ast}}}1
    {•}{{\ensuremath{\bullet}}}1,
  % Comments
  morecomment=[l]{--},
  morecomment=[s]{/-}{-/},
  morestring=[b]",
  stringstyle=\color{stringcolor},
  keepspaces=true,
  keywordstyle=[1]\color{keywordcolor}\bfseries,
  keywordstyle=[2]\color{sortcolor}\bfseries,
  keywordstyle=[3]\color{tacticcolor}\bfseries,
  keywordstyle=[4]\color{errorcolor}\bfseries\underline,
  escapeinside={(*@}{@*)},
  columns=fullflexible,
  basicstyle=\ttfamily\small,
  commentstyle=\color{commentcolor}\itshape,
  breaklines=true,
  showstringspaces=false,
  frame=single,
  framesep=2mm,
  xleftmargin=4mm,
  numbers=left,
  numberstyle=\tiny\color{commentcolor},
  numbersep=6pt,
}

% Inline Lean code in prose: \lean{Int.ModEq.pow}
\newcommand{\lean}[1]{\lstinline[language=lean,basicstyle=\ttfamily\normalsize]|#1|}

\lstset{language=lean}

% --- XeTeX rendering patch (2026-10-06) --- tectonic/XeTeX does not apply
% the listings literate table to multi-byte characters, so every non-ASCII,
% non-CJK codepoint that occurs inside THIS paper's listings is mapped to a
% math glyph as an active character via the standard \lccode/\lowercase idiom
% (ASCII-only source, no extra packages). The mapped set is verified against
% the spliced blocks by audit/build.py at assembly time -- an unmapped
% codepoint aborts the build. Deliberately NOT mapped: U+00B7, U+2014, U+2026
% and all CJK ideographs / CJK punctuation --- they route through xeCJK +
% SimSun (making them active first aborts the run, observed in earlier
% packages). Rendering only: listing source bytes are untouched, so the
% byte-exactness checks are unaffected, and the \ifdefined guard leaves the
% pdflatex path identical to the template.
\ifdefined\XeTeXversion
  {\lccode`\~="00A7 \lowercase{\global\catcode`~=\active \global\def~{\S}}}
  {\lccode`\~="00AC \lowercase{\global\catcode`~=\active \global\def~{\ensuremath{\neg}}}}
  {\lccode`\~="00B2 \lowercase{\global\catcode`~=\active \global\def~{\ensuremath{{}^2}}}}
  {\lccode`\~="03BD \lowercase{\global\catcode`~=\active \global\def~{\ensuremath{\nu}}}}
  {\lccode`\~="2115 \lowercase{\global\catcode`~=\active \global\def~{\ensuremath{\mathbb{N}}}}}
  {\lccode`\~="2192 \lowercase{\global\catcode`~=\active \global\def~{\ensuremath{\rightarrow}}}}
  {\lccode`\~="2194 \lowercase{\global\catcode`~=\active \global\def~{\ensuremath{\leftrightarrow}}}}
  {\lccode`\~="2200 \lowercase{\global\catcode`~=\active \global\def~{\ensuremath{\forall}}}}
  {\lccode`\~="2203 \lowercase{\global\catcode`~=\active \global\def~{\ensuremath{\exists}}}}
  {\lccode`\~="2208 \lowercase{\global\catcode`~=\active \global\def~{\ensuremath{\in}}}}
  {\lccode`\~="2211 \lowercase{\global\catcode`~=\active \global\def~{\ensuremath{\sum}}}}
  {\lccode`\~="2212 \lowercase{\global\catcode`~=\active \global\def~{\ensuremath{-}}}}
  {\lccode`\~="2223 \lowercase{\global\catcode`~=\active \global\def~{\ensuremath{\mid}}}}
  {\lccode`\~="2224 \lowercase{\global\catcode`~=\active \global\def~{\ensuremath{\nmid}}}}
  {\lccode`\~="2227 \lowercase{\global\catcode`~=\active \global\def~{\ensuremath{\wedge}}}}
  {\lccode`\~="2260 \lowercase{\global\catcode`~=\active \global\def~{\ensuremath{\neq}}}}
  {\lccode`\~="2264 \lowercase{\global\catcode`~=\active \global\def~{\ensuremath{\leq}}}}
  {\lccode`\~="2265 \lowercase{\global\catcode`~=\active \global\def~{\ensuremath{\geq}}}}
  {\lccode`\~="2270 \lowercase{\global\catcode`~=\active \global\def~{\ensuremath{\nleq}}}}
  {\lccode`\~="230A \lowercase{\global\catcode`~=\active \global\def~{\ensuremath{\lfloor}}}}
  {\lccode`\~="230B \lowercase{\global\catcode`~=\active \global\def~{\ensuremath{\rfloor}}}}
  {\lccode`\~="27E8 \lowercase{\global\catcode`~=\active \global\def~{\ensuremath{\langle}}}}
  {\lccode`\~="27E9 \lowercase{\global\catcode`~=\active \global\def~{\ensuremath{\rangle}}}}
  {\lccode`\~="27F8 \lowercase{\global\catcode`~=\active \global\def~{\ensuremath{\Longleftarrow}}}}
  {\lccode`\~="27F9 \lowercase{\global\catcode`~=\active \global\def~{\ensuremath{\Longrightarrow}}}}
  {\lccode`\~="27FA \lowercase{\global\catcode`~=\active \global\def~{\ensuremath{\Longleftrightarrow}}}}
\fi
\ifdefined\XeTeXversion
  \usepackage{xeCJK}
  \setCJKmainfont{SimSun}
\fi

\newtheorem{theorem}{Theorem}
\newtheorem{lemma}[theorem]{Lemma}
\newtheorem{corollary}[theorem]{Corollary}
\theoremstyle{definition}
\newtheorem{definition}[theorem]{Definition}

\title{The Diophantine equation $x^2 = n! \cdot n\#$: the complete solution
set,\\ formalized and machine-checked in Lean 4}
\author{Aurora0134\thanks{Contact: via the GitHub account Aurora0134.}}
\date{October 6, 2026}

\begin{document}
\maketitle

\begin{abstract}
We determine the complete solution set of the Diophantine equation
$x^2 = n! \cdot n\#$ (the factorial times the primorial): for
$n \in \mathbb{N}$ the equation is solvable exactly for
$n \in \{0,1,2,3,4,5\}$, with $x = 1, 1, 2, 6, 12, 60$. This settles the
first of the two cases that Novakovi\'c (2026) recorded as believed-finite
with the words ``at the moment we were not able to give an argument''. The
proof combines a valuation criterion ($n! \cdot n\#$ is a square if and only
if $\nu_p(n!)$ is odd for every prime $p \le n$), an interval-prime lemma (for
$n \ge 22$ there is a prime in $(n/3, n/2]$) proved by a prime-certificate
chain up to $10^{6}$ and an Erd\H{o}s-style factorial-ratio estimate beyond
it, and explicit witnesses for small $n$. Every statement is machine-checked
in Lean~4 against mathlib v4.34.0: the statements were frozen by hash before
proving, the proof compiles with no \texttt{sorry}, and
\texttt{\#print axioms} reports only \{propext, \texttt{Quot.sound},
\texttt{Classical.choice}\}. The results were developed with AI assistance
under human direction, and the Lean 4 kernel is the sole adjudicator. Lean
sources, frozen statements, and audit logs accompany the paper as a deposited
artifact, and a companion claim of record is deposited separately.
\end{abstract}

\section{Introduction}\label{sec:intro}

The equation of the title asks when the product of the factorial and the
primorial is a perfect square. It appears, together with the companion
equation $x^2 = n!! \cdot [1, \ldots, n]$, in Novakovi\'c's 2026 study
\cite{novakovic2026} of the Diophantine equation
$P(x) = \prod_{i=1}^{r} H_{n_i}$, where the author writes (verbatim wording;
mathematical notation retypeset): ``We believe that finitely many solutions
occur for $x^2 = n! \cdot n\#$ or $x^2 = n!! \cdot [1, \ldots, n]$, but at
the moment we were not able to give an argument''. The present paper answers
the first of these two cases completely: the equation
$x^2 = n! \cdot n\#$ has, for $n \in \mathbb{N}$, exactly the solutions
$n \in \{0,1,2,3,4,5\}$. The second case is not claimed here.

Three contributions are recorded together. First, the mathematical result
itself, as an if-and-only-if characterization with explicit witnesses
(Section~\ref{sec:results}); an exclusivity review found no
occupying record for this statement on the channels it could reach
(Section~\ref{sec:related}). Second, a complete machine-checked formalization:
eight statements frozen by hash before proving, proved in Lean~4 against
mathlib with no placeholder and with axiom dependencies exactly the three
standard ones (Section~\ref{sec:verify}). Third, the artifact is built so
that a reader can re-run the adjudication --- the frozen statement's hash,
the strict compile's exit status and the axiom printout all ship with the
deposit --- rather than take the text's word for it. This is why the paper
addresses a formal-methods readership as well as a number-theoretic one.

A companion claim of record for the same kernel result was deposited
separately \cite{claimdeposit}; the two deposits
are byte-level siblings at the level of the proof sources and differ in
genre: the claim note is a priority record, this paper is the narrative
account.

\subsection*{Notation}

We write $n!$ for the factorial and $n\# = \prod_{p \le n,\ p \text{ prime}} p$
for the primorial; $\nu_p(m)$ denotes the exponent of the prime $p$ in $m$.
In the formalization, \lean{primorial} is the top-level mathlib declaration
and $\nu_p(m)$ is \lean{Nat.factorization}. Division in the interval-prime
lemma (Lemma~\ref{lem:B}) is natural-number division, so the interval reads
$\lfloor n/3 \rfloor < p \le \lfloor n/2 \rfloor$; for prime existence this
coincides with the real reading, since $p$ is an integer.

\section{Mathematical background}\label{sec:background}

\begin{definition}\label{def:primorial}
For $n \in \mathbb{N}$, the \emph{primorial} $n\#$ is the product of the
primes $p \le n$ (so $0\# = 1\# = 1$, $5\# = 30$). For a prime $p$ and
$m \ge 1$, $\nu_p(m)$ is the exponent of $p$ in the prime factorization of
$m$.
\end{definition}

The only classical ingredient the proof uses beyond arithmetic is Legendre's
formula $\nu_p(n!) = \sum_{i \ge 1} \lfloor n/p^{i} \rfloor$, which enters
through mathlib's \lean{Nat.factorization_factorial}. The shape of the
argument is: a square has all valuations even (Lemma~\ref{lem:A} reduces the
equation to a parity condition on $\nu_p(n!)$); a prime $p$ with
$2p \le n < 3p$ contributes $\nu_p(n!) = \lfloor n/p \rfloor = 2$, an even
number, provided $p^2 > n$ kills the higher terms; and such a prime exists
for all $n \ge 22$ (Lemma~\ref{lem:B}). The small cases are decided
pointwise.

\paragraph{Where this sits in the literature.} The equation belongs to the
family around the Brocard--Ramanujan equation $n! + 1 = m^2$, of which
Novakovi\'c's equation $P(x) = \prod_{i=1}^{r} H_{n_i}$
\cite{novakovic2026} is a generalization; the same author's companion paper
\cite{novakovic2026b} is recorded here with the restricted wording that its
object of study does not contain $n\#$. Adjacent objects and records that
must be distinguished from the present statement --- the squarefree part of
$n!$ (OEIS A055204 \cite{oeisA055204}), the disproved Erd\H{o}s problem
\#399 \cite{erdos399}, primorial-core Pell equations \cite{peixoto2026},
square products of factorials \cite{yudin2026}, and the interval-prime
literature (Nagura \cite{nagura1952}, OEIS A166968 \cite{oeisA166968}) ---
are delimited one by one in Section~\ref{sec:related}.

\section{Formalisation setup}\label{sec:setup}

The formalization is written against Lean~4 \cite{lean4}, version
\texttt{v4.34.0} (the \texttt{leanprover/lean4} release of that version),
with mathlib4 \cite{mathlib} \texttt{v4.34.0} at revision \texttt{5ed29652};
the pin ships with the artifact (\texttt{proofs/lean-toolchain}). The eight
statements were frozen by hash before proving began; the proving stage may
replace proof bodies only, and the final audit re-compares the delivered
statement headers against the frozen snapshot under two whitespace
conventions.

\paragraph{Encoding decisions.} Four encoding choices are worth recording.
(i) The primorial is mathlib's top-level \lean{primorial}; the primorial
valuations enter through the squarefreeness of \lean{primorial} and the
divisibility characterization \lean{Nat.Prime.dvd_primorial_iff}
(\lean{primorial_factorization_eq_one}, \lean{primorial_factorization_eq_zero}).
(ii) The main result is delivered in two forms: the iff form
\lean{isSquare_factorial_mul_primorial_iff} (``$n! \cdot n\#$ is a square if
and only if $n \le 5$'') and the existence-witness form
\lean{exists_sq_factorial_mul_primorial_iff} (``$x^2 = n! \cdot n\#$ is
solvable if and only if $n \in \{0,1,2,3,4,5\}$''), bridged by
\lean{isSquare_iff_exists_sq}; on $\mathbb{N}$, $n \le 5$ and
$n \in \{0,\ldots,5\}$ coincide. (iii) The interval-prime lemma is stated
with natural-number division; this is faithful to the real reading for the
purpose of prime existence, as recorded in the formalization's fidelity
note. (iv) The criterion of Lemma~\ref{lem:A} carries the hypothesis
$1 \le n$ as frozen; the proof does not use it (the criterion holds at
$n = 0$ as well), and the lemma is stated as frozen rather than strengthened
--- Section~\ref{sec:limits}(e).

\paragraph{Audit mechanism.} The statement is frozen by hash; a strict
compile must exit $0$ with no placeholder (\texttt{sorry}/\texttt{admit})
anywhere in the file, comments included; \texttt{\#print axioms} must report
only the three standard axioms; and the statement comparison is re-run by an
independent audit at the end. The verdicts for this paper are reproduced in
Section~\ref{sec:verify} and ship with the artifact.

\section{Main results}\label{sec:results}

% LEAN: tasks/20261006-nfact-primorial-square :: isSquare_factorial_mul_primorial_iff grade=完全证明
\begin{theorem}[complete solution set, iff form]\label{thm:main}
For every $n \in \mathbb{N}$, $n! \cdot n\#$ is a perfect square if and only
if $n \le 5$.
\end{theorem}

% LEAN: tasks/20261006-nfact-primorial-square :: exists_sq_factorial_mul_primorial_iff grade=完全证明
\begin{theorem}[existence-witness form]\label{thm:exists}
For every $n \in \mathbb{N}$, the equation $x^2 = n! \cdot n\#$ has a
solution $x \in \mathbb{N}$ if and only if $n \in \{0,1,2,3,4,5\}$.
\end{theorem}

% LEAN: tasks/20261006-nfact-primorial-square :: isSquare_iff_all_factorization_odd grade=完全证明
\begin{lemma}[valuation criterion]\label{lem:A}
For $n \ge 1$: $n! \cdot n\#$ is a perfect square if and only if
$\nu_p(n!)$ is odd for every prime $p \le n$.
\end{lemma}

% LEAN: tasks/20261006-nfact-primorial-square :: exists_prime_between_div_three_and_div_two grade=完全证明
\begin{lemma}[interval-prime lemma]\label{lem:B}
For every integer $n \ge 22$ there exists a prime $p$ with
$n/3 < p \le n/2$.
\end{lemma}

% LEAN: tasks/20261006-nfact-primorial-square :: witnesses_zero_to_five grade=完全证明
\begin{theorem}[solution witnesses]\label{thm:witnesses}
$1^2 = 0! \cdot 0\#$, $1^2 = 1! \cdot 1\#$, $2^2 = 2! \cdot 2\#$,
$6^2 = 3! \cdot 3\#$, $12^2 = 4! \cdot 4\#$, and $60^2 = 5! \cdot 5\#$.
\end{theorem}

\noindent Three auxiliary factorization statements belong to the same frozen
set and carry the same grade (verbatim in Listing~\ref{lst:stmt}):
$\nu_p(n\#) = 1$ for primes $p \le n$, $\nu_p(n\#) = 0$ for primes
$p \nleq n$, and $\nu_p(n! \cdot n\#) = \nu_p(n!) + 1$ for primes $p \le n$.
% LEAN: tasks/20261006-nfact-primorial-square :: primorial_factorization_eq_one grade=完全证明
% LEAN: tasks/20261006-nfact-primorial-square :: primorial_factorization_eq_zero grade=完全证明
% LEAN: tasks/20261006-nfact-primorial-square :: factorization_factorial_mul_primorial grade=完全证明

\begin{lstlisting}[caption={Frozen statement snapshot (Lean 4): the complete
frozen set, copied from the \texttt{lean} fence of the snapshot shipped as
\texttt{proofs/nfact-primorial-square.statement.txt} (a comment-cleaned
edition: docstrings translated to English for publication; the statement
text is byte-identical to the frozen snapshot of record, sha256
\texttt{0b13d08d}; correspondence record
\texttt{audit/comment-cleanup-diff.md}).},label={lst:stmt}]
/-- **Main theorem (complete solution set, iff form)**:
    `n! * n#` is a perfect square ⟺ `n ≤ 5` (for `n : ℕ`).
    Over ℕ, `n ≤ 5 ⟺ n ∈ {0,1,2,3,4,5}`, so this is the main statement "the solutions are exactly n ∈ {0,...,5}". -/
theorem isSquare_factorial_mul_primorial_iff (n : ℕ) :
    IsSquare (n ! * primorial n) ↔ n ≤ 5

/-- **Main theorem (existence-witness form, delivered alongside the iff form)**:
    the Diophantine equation `x² = n! * n#` has a solution ⟺ `n ∈ {0,1,2,3,4,5}`. -/
theorem exists_sq_factorial_mul_primorial_iff (n : ℕ) :
    (∃ x : ℕ, x ^ 2 = n ! * primorial n) ↔ n ∈ ({0, 1, 2, 3, 4, 5} : Finset ℕ)

/-- **Lemma A (equivalent valuation criterion)**: for `n ≥ 1`, `n! * n#` is a perfect square ⟺
    `ν_p(n!)` is odd for every prime `p ≤ n`. -/
theorem isSquare_iff_all_factorization_odd (n : ℕ) (hn : 1 ≤ n) :
    IsSquare (n ! * primorial n) ↔
      ∀ p : ℕ, p.Prime → p ≤ n → Odd ((n !).factorization p)

/-- **Lemma B (interval prime lemma)**: for every integer `n ≥ 22` there exists a prime `p` with `n/3 < p ≤ n/2`. -/
theorem exists_prime_between_div_three_and_div_two (n : ℕ) (hn : 22 ≤ n) :
    ∃ p : ℕ, p.Prime ∧ n / 3 < p ∧ p ≤ n / 2

/-- **Witness layer (explicit square roots for the ⟸ direction of the main theorem)**: `x² = n! * n#` at each of n = 0..5, with x = 1, 1, 2, 6, 12, 60. -/
theorem witnesses_zero_to_five :
    1 ^ 2 = 0 ! * primorial 0 ∧
    1 ^ 2 = 1 ! * primorial 1 ∧
    2 ^ 2 = 2 ! * primorial 2 ∧
    6 ^ 2 = 3 ! * primorial 3 ∧
    12 ^ 2 = 4 ! * primorial 4 ∧
    60 ^ 2 = 5 ! * primorial 5

/-- Auxiliary statement: `ν_p(n#) = 1` (primes `p ≤ n`). -/
theorem primorial_factorization_eq_one {p n : ℕ} (hp : p.Prime) (hpn : p ≤ n) :
    (primorial n).factorization p = 1

/-- Auxiliary statement: `ν_p(n#) = 0` (primes `p ≰ n`). -/
theorem primorial_factorization_eq_zero {p n : ℕ} (hp : p.Prime) (hpn : ¬ p ≤ n) :
    (primorial n).factorization p = 0

/-- Auxiliary statement (the multiplicative skeleton of the criterion): `ν_p(n! * n#) = ν_p(n!) + 1` (primes `p ≤ n`). -/
theorem factorization_factorial_mul_primorial {p n : ℕ} (hp : p.Prime) (hpn : p ≤ n) :
    (n ! * primorial n).factorization p = (n !).factorization p + 1
\end{lstlisting}

\section{Proofs}\label{sec:proofs}

The final proof file (1094 lines) is a single Lean module in two
namespaces, and the proof assembles four components: the valuation criterion,
the exclusion of large $n$ via the interval-prime lemma, the small-$n$
witness layer, and the assembly of the two main theorems. The file ships
byte-exactly in the artifact under \texttt{proofs/}; the excerpts below are
byte-identical contiguous blocks of it (independent mechanical record:
\texttt{audit/listing-verify.txt}, which ships in the deposit). The shipped
file is the comment-cleaned edition (English docstrings); its code is
byte-identical to the frozen development (record:
\texttt{audit/comment-cleanup-diff.md}). We describe
the proof in the order in which it is built.

\subsection{The valuation criterion}

The multiplicative spine of the argument is
\lean{factorization_factorial_mul_primorial}: for a prime $p \le n$, the
primorial contributes $p$ exactly once (it is squarefree and divisible by
$p$), so
\[
\nu_p(n! \cdot n\#) = \nu_p(n!) + 1.
\]
A natural number is a square if and only if all its prime valuations are
even (mathlib's \lean{Nat.isSquare_iff_even_factorization}). For $p \le n$
the displayed identity makes the evenness of $\nu_p(n! \cdot n\#)$ equivalent
to the oddness of $\nu_p(n!)$; for $p > n$ the prime divides neither $n!$
nor $n\#$, so the valuation is $0$, an even number. This proves
Lemma~\ref{lem:A} in both directions. The backward direction is discharged
through the library bridge rather than by an explicit square-root
construction --- a design-stage rerouting (evidence by kernel pre-checks),
not a weakening.

\subsection{The witness layer}

The six solutions are exhibited pointwise (Theorem~\ref{thm:witnesses},
closed by decision). For the remaining small cases, the range
$6 \le n \le 21$ is excluded by sixteen explicit factorizations: each point
exhibits a prime $p$ with $\nu_p(n! \cdot n\#) = 2e + 1$ odd --- for
instance $\nu_2(6! \cdot 6\#) = 5$ since $6! \cdot 6\# = 2^{5} \cdot 675$
with $2 \nmid 675$ --- and the bridge
\lean{not_isSquare_of_odd_factorization} turns an odd valuation into a
non-square (a square $m = r \cdot r$ would have
$\nu_p(m) = \nu_p(r) + \nu_p(r)$ even). The assembly
\lean{nonsquare_six_to_twenty_one} closes the range by case analysis. The
sixteen witnesses appear in full in Listing~\ref{lst:main}.

\subsection{The interval-prime lemma: the certificate chain}

For $22 \le n \le 999999$ the lemma is closed by a certificate chain
(Listing~\ref{lst:chain}). The single hop \lean{lemmaB_hop} observes that a
prime $q$ with $2q \le n \le 3q - 1$ witnesses $n/3 < q \le n/2$. A greedy
chain of twenty-nine primes,
\[
11, 13, 19, 23, 31, 43, 61, 89, 131, 193, 283, 421, 631, 941, 1409, 2113,
3169, 4751, 7121, 10667, \ldots, 409609,
\]
together with twenty-eight continuity certificates $2q' \le 3q$ between
consecutive primes, covers the segment from $22 = 2 \cdot 11$ up to
$3 \cdot 409609 - 1 = 1228826 \ge 999999$; every certificate is a one-line
\texttt{norm\_num}/\texttt{omega} fact, and the chain's decision tree is
\lean{lemmaB_upto_999999}. The chain was produced by a mechanical generator
and kernel-checked as a whole; the segment $22 \le n \le 32$ is covered by
the first hop with the uniform witness $p = 11$.

\begin{lstlisting}[basicstyle=\ttfamily\footnotesize,caption={Excerpt of the
proof file (lines 96--168 of \texttt{final-nfact-primorial-square.lean},
which ships in full in \texttt{proofs/}): the certificate-chain mechanism for
Lemma~\ref{lem:B} on $22 \le n \le 999999$ --- the two boundary certificates,
the single-hop lemma \texttt{lemmaB\_hop}, and the twenty-nine-hop decision
tree.},label={lst:chain}]
theorem chain_start : 22 = 2 * 11 := by omega
theorem chain_end : 999999 ≤ 3 * 409609 - 1 := by omega

/-- Certificate-chain single hop: for prime `q` and `n ∈ [2q, 3q - 1]`,
    `q` itself witnesses `n/3 < q ≤ n/2`.
    Note: extract `hq.two_le` first, otherwise `omega` gets stuck at the `q = 0` counterexample under saturating subtraction on ℕ. -/
theorem lemmaB_hop {q : ℕ} (hq : q.Prime) {n : ℕ} (h1 : 2 * q ≤ n) (h2 : n ≤ 3 * q - 1) :
    ∃ p : ℕ, p.Prime ∧ n / 3 < p ∧ p ≤ n / 2 :=
  have hq2 : 2 ≤ q := hq.two_le
  ⟨q, hq, by omega, by omega⟩

/-- **Lemma B, finite segment (certificate-chain decision tree)**: for every `n ∈ [22, 999999]`,
    `∃ p, p.Prime ∧ n/3 < p ∧ p ≤ n/2`.
    Joins the tail theorem (which starts at n = 1000000) by a case split on `n < 1000000`. -/
theorem lemmaB_upto_999999 (n : ℕ) (h1 : 22 ≤ n) (h2 : n ≤ 999999) :
    ∃ p : ℕ, p.Prime ∧ n / 3 < p ∧ p ≤ n / 2 := by
  rcases lt_or_ge n 33 with h | h
  · exact lemmaB_hop (q := 11) prime_cert_11 (by omega) (by omega)
  · rcases lt_or_ge n 39 with h | h
    · exact lemmaB_hop (q := 13) prime_cert_13 (by omega) (by omega)
    · rcases lt_or_ge n 57 with h | h
      · exact lemmaB_hop (q := 19) prime_cert_19 (by omega) (by omega)
      · rcases lt_or_ge n 69 with h | h
        · exact lemmaB_hop (q := 23) prime_cert_23 (by omega) (by omega)
        · rcases lt_or_ge n 93 with h | h
          · exact lemmaB_hop (q := 31) prime_cert_31 (by omega) (by omega)
          · rcases lt_or_ge n 129 with h | h
            · exact lemmaB_hop (q := 43) prime_cert_43 (by omega) (by omega)
            · rcases lt_or_ge n 183 with h | h
              · exact lemmaB_hop (q := 61) prime_cert_61 (by omega) (by omega)
              · rcases lt_or_ge n 267 with h | h
                · exact lemmaB_hop (q := 89) prime_cert_89 (by omega) (by omega)
                · rcases lt_or_ge n 393 with h | h
                  · exact lemmaB_hop (q := 131) prime_cert_131 (by omega) (by omega)
                  · rcases lt_or_ge n 579 with h | h
                    · exact lemmaB_hop (q := 193) prime_cert_193 (by omega) (by omega)
                    · rcases lt_or_ge n 849 with h | h
                      · exact lemmaB_hop (q := 283) prime_cert_283 (by omega) (by omega)
                      · rcases lt_or_ge n 1263 with h | h
                        · exact lemmaB_hop (q := 421) prime_cert_421 (by omega) (by omega)
                        · rcases lt_or_ge n 1893 with h | h
                          · exact lemmaB_hop (q := 631) prime_cert_631 (by omega) (by omega)
                          · rcases lt_or_ge n 2823 with h | h
                            · exact lemmaB_hop (q := 941) prime_cert_941 (by omega) (by omega)
                            · rcases lt_or_ge n 4227 with h | h
                              · exact lemmaB_hop (q := 1409) prime_cert_1409 (by omega) (by omega)
                              · rcases lt_or_ge n 6339 with h | h
                                · exact lemmaB_hop (q := 2113) prime_cert_2113 (by omega) (by omega)
                                · rcases lt_or_ge n 9507 with h | h
                                  · exact lemmaB_hop (q := 3169) prime_cert_3169 (by omega) (by omega)
                                  · rcases lt_or_ge n 14253 with h | h
                                    · exact lemmaB_hop (q := 4751) prime_cert_4751 (by omega) (by omega)
                                    · rcases lt_or_ge n 21363 with h | h
                                      · exact lemmaB_hop (q := 7121) prime_cert_7121 (by omega) (by omega)
                                      · rcases lt_or_ge n 32001 with h | h
                                        · exact lemmaB_hop (q := 10667) prime_cert_10667 (by omega) (by omega)
                                        · rcases lt_or_ge n 47973 with h | h
                                          · exact lemmaB_hop (q := 15991) prime_cert_15991 (by omega) (by omega)
                                          · rcases lt_or_ge n 71943 with h | h
                                            · exact lemmaB_hop (q := 23981) prime_cert_23981 (by omega) (by omega)
                                            · rcases lt_or_ge n 107907 with h | h
                                              · exact lemmaB_hop (q := 35969) prime_cert_35969 (by omega) (by omega)
                                              · rcases lt_or_ge n 161853 with h | h
                                                · exact lemmaB_hop (q := 53951) prime_cert_53951 (by omega) (by omega)
                                                · rcases lt_or_ge n 242769 with h | h
                                                  · exact lemmaB_hop (q := 80923) prime_cert_80923 (by omega) (by omega)
                                                  · rcases lt_or_ge n 364137 with h | h
                                                    · exact lemmaB_hop (q := 121379) prime_cert_121379 (by omega) (by omega)
                                                    · rcases lt_or_ge n 546177 with h | h
                                                      · exact lemmaB_hop (q := 182059) prime_cert_182059 (by omega) (by omega)
                                                      · rcases lt_or_ge n 819249 with h | h
                                                        · exact lemmaB_hop (q := 273083) prime_cert_273083 (by omega) (by omega)
                                                        · exact lemmaB_hop (q := 409609) prime_cert_409609 (by omega) (by omega)
\end{lstlisting}

\subsection{The interval-prime lemma: the factorial-ratio tail}

For $n \ge 10^{6}$ the lemma is closed by an Erd\H{o}s-style factorial-ratio
argument (Listing~\ref{lst:tail}). With
\[
\mathrm{tailR}(n) = \frac{(n/2)!\,(n/6)!}{(n/3)!\,(n/4)!\,(n/12)!}
\]
(natural-number division throughout), the valuations are
$\nu_p(\mathrm{tailR}(n)) = \sum_{i \ge 1} \mathrm{hval}(\lfloor n/p^{i}
\rfloor)$, where $\mathrm{hval}(m) = \lfloor m/2 \rfloor + \lfloor m/6
\rfloor - \lfloor m/3 \rfloor - \lfloor m/4 \rfloor - \lfloor m/12 \rfloor$
takes only the values $0$ and $1$, with $\mathrm{hval}(m) = 1$ exactly when
$m \bmod 12 \in \{2,6,7,8,10,11\}$. Two bounds collide. On the lower side, an
induction on blocks of twelve (\lean{tailR_lower_twelve},
\lean{tailR_lower_strict}) gives the strict estimate
\[
18^{\lfloor n/12 \rfloor} < \mathrm{tailR}(n) \cdot (4\lfloor n/12
\rfloor + 3)^{6}.
\]
On the upper side, if the interval $(n/3, n/2]$ contained no prime, then the
positions with $\mathrm{hval} = 1$ in the critical range would be excluded,
and the total mass collapses to
$\mathrm{tailR}(n) \le (n/2)^{\sqrt{n/2}} \cdot 4^{n/6}$
(\lean{tailR_le_of_no_prime}), with $4^{n/6} \le 4 \cdot 16^{\lfloor n/12
\rfloor}$. The decisive closure
\[
4 \cdot (4q+3)^{6} \cdot (6q+5)^{\sqrt{6q+5}} \cdot 16^{q} \le 18^{q}
\qquad (q \ge 10^{5})
\]
(\lean{tail_main_inequality_nat}) is a real-variable concavity two-point
estimate (\lean{concave_log_affine}, \lean{real_tail_inequality}), modelled
on mathlib's Bertrand proof; the margin is the bare inequality $18 > 16$.
The join at the boundary, $10^{6} \le n < 1.2 \cdot 10^{6}$, is a single
certificate hop with $q = 400009$.

\begin{lstlisting}[basicstyle=\ttfamily\footnotesize,caption={Excerpt of the
proof file (lines 735--767): the tail theorem for $n \ge 10^{6}$
(\texttt{lemmaB\_tail\_1000000}), closing Lemma~\ref{lem:B} on the infinite
regime via the factorial-ratio estimate.},label={lst:tail}]
theorem lemmaB_tail_1000000 (n : ℕ) (hn : 1000000 ≤ n) :
    ∃ p : ℕ, p.Prime ∧ n / 3 < p ∧ p ≤ n / 2 := by
  rcases lt_or_ge n 1200000 with h | h
  · exact lemmaB_hop (q := 400009) (by norm_num) (by omega) (by omega)
  · by_contra hc
    push_neg at hc
    have hc' : ∀ p : ℕ, p.Prime → n / 3 < p → ¬ p ≤ n / 2 := by
      intro p hp1 hp2 hle
      have hcc := hc p hp1 hp2
      omega
    have hub := tailR_le_of_no_prime n (by omega) hc'
    have hlb := tailR_lower_strict n hn
    have hmain := tail_main_inequality_nat (n/12) (by omega)
    have h2 : n < (n/12 + 1) * 12 := (Nat.div_lt_iff_lt_mul (by norm_num : (0:ℕ) < 12)).mp (Nat.lt_succ_self _)
    have e1 : (n/2)^(Nat.sqrt (n/2)) ≤ (6*(n/12)+5)^(Nat.sqrt (6*(n/12)+5)) := by
      have hbase : n/2 ≤ 6*(n/12)+5 := by omega
      have hs : Nat.sqrt (n/2) ≤ Nat.sqrt (6*(n/12)+5) := Nat.sqrt_le_sqrt hbase
      exact (Nat.pow_le_pow_right (by omega : 0 < n/2) hs).trans (Nat.pow_le_pow_left hbase _)
    have e2 : 4^(n/6) ≤ 4 * 16^(n/12) := by
      have h4 : n/6 ≤ 2*(n/12) + 1 := by omega
      calc 4^(n/6) ≤ 4^(2*(n/12)+1) := Nat.pow_le_pow_right (by norm_num) h4
        _ = 4 * 16^(n/12) := by
          rw [pow_succ, pow_mul]
          norm_num
          ring
    have key : 18^(n/12) < 18^(n/12) := by
      calc 18^(n/12) < tailR n * (4*(n/12)+3)^6 := hlb
        _ ≤ ((n/2)^(Nat.sqrt (n/2)) * 4^(n/6)) * (4*(n/12)+3)^6 := Nat.mul_le_mul_right _ hub
        _ ≤ ((6*(n/12)+5)^(Nat.sqrt (6*(n/12)+5)) * (4 * 16^(n/12))) * (4*(n/12)+3)^6 :=
            Nat.mul_le_mul_right _ (Nat.mul_le_mul e1 e2)
        _ = 4 * (4*(n/12)+3)^6 * (6*(n/12)+5)^(Nat.sqrt (6*(n/12)+5)) * 16^(n/12) := by ring
        _ ≤ 18^(n/12) := hmain
    exact absurd key (lt_irrefl _)
\end{lstlisting}

\subsection{Assembly of the main theorems}

The backward direction of Theorem~\ref{thm:main} is the witness layer: for
$n \le 5$, case analysis (\texttt{interval\_cases}) produces the explicit
square roots $x = 1, 1, 2, 6, 12, 60$ via \lean{isSquare_iff_exists_sq}. For
the forward direction, suppose $n! \cdot n\#$ is a square and $n \ge 6$. If
$6 \le n \le 21$, the witness layer contradicts it directly. If $n \ge 22$,
Lemma~\ref{lem:B} yields a prime $p$ with $n/3 < p \le n/2$; then
$2p \le n < 3p$ and $p^{2} \ge 3p > n$ (using $p \ge 3$), so Legendre's
formula collapses to its first term:
\[
\nu_p(n!) = \sum_{i \ge 1} \lfloor n/p^{i} \rfloor = \lfloor n/p \rfloor = 2
\]
(in the formalization, \lean{Nat.factorization_factorial} is applied with
the bound $b = 2$, since $p^{2} > n$ gives \lean{Nat.log} $p\ n < 2$ via
\lean{Nat.log_lt_of_lt_pow}). But the criterion of Lemma~\ref{lem:A} forces
$\nu_p(n!)$ to be odd --- contradiction. Theorem~\ref{thm:exists} is then the
bridge through \lean{isSquare_iff_exists_sq} together with the finite-set
reading of $n \le 5$ on $\mathbb{N}$.

\begin{lstlisting}[basicstyle=\ttfamily\footnotesize,caption={Excerpt of the
proof file (lines 771--1094): the main namespace \texttt{NFactPrimorialSquare}
in full --- the three auxiliary factorization lemmas, Lemma~\ref{lem:A}, the
witness layer, the sixteen non-square witnesses, Lemma~\ref{lem:B} assembled
from the chain and the tail, and the two main theorems.},label={lst:main}]
namespace NFactPrimorialSquare

/-- Auxiliary statement: `ν_p(n#) = 1` for primes `p ≤ n`. -/
theorem primorial_factorization_eq_one {p n : ℕ} (hp : p.Prime) (hpn : p ≤ n) :
    (primorial n).factorization p = 1 :=
  factorization_eq_one_of_squarefree (squarefree_primorial n) hp
    ((Nat.Prime.dvd_primorial_iff hp).mpr hpn)

/-- Auxiliary statement: `ν_p(n#) = 0` for primes `p ≰ n`. -/
theorem primorial_factorization_eq_zero {p n : ℕ} (hp : p.Prime) (hpn : ¬ p ≤ n) :
    (primorial n).factorization p = 0 :=
  Nat.factorization_eq_zero_of_not_dvd
    (fun h => hpn ((Nat.Prime.dvd_primorial_iff hp).mp h))

/-- Auxiliary statement (the multiplicative skeleton of the criterion): `ν_p(n! * n#) = ν_p(n!) + 1` for primes `p ≤ n`. -/
theorem factorization_factorial_mul_primorial {p n : ℕ} (hp : p.Prime) (hpn : p ≤ n) :
    (n ! * primorial n).factorization p = (n !).factorization p + 1 := by
  rw [Nat.factorization_mul (Nat.factorial_ne_zero n) (primorial_ne_zero n),
      Finsupp.add_apply, primorial_factorization_eq_one hp hpn]

/-- **Lemma A (equivalent valuation criterion)**: for `n ≥ 1`, `n! * n#` is a perfect square ⟺
    `ν_p(n!)` is odd for every prime `p ≤ n`.
    (`Nat.factorization m p` is ν_p(m): the exponent of the prime `p` in `m`.) -/
theorem isSquare_iff_all_factorization_odd (n : ℕ) (hn : 1 ≤ n) :
    IsSquare (n ! * primorial n) ↔
      ∀ p : ℕ, p.Prime → p ≤ n → Odd ((n !).factorization p) := by
  constructor
  · intro h p hp hpn
    have hEven : Even ((n ! * primorial n).factorization p) :=
      Nat.isSquare_iff_even_factorization.mp h p hp
    rw [factorization_factorial_mul_primorial hp hpn] at hEven
    obtain ⟨m, hm⟩ := hEven
    cases m with
    | zero => simp at hm
    | succ k => exact ⟨k, by omega⟩
  · intro h
    apply Nat.isSquare_iff_even_factorization.mpr
    intro p hp
    by_cases hpn : p ≤ n
    · rw [factorization_factorial_mul_primorial hp hpn]
      exact (h p hp hpn).add_one
    · rw [Nat.factorization_mul (Nat.factorial_ne_zero n) (primorial_ne_zero n),
        Finsupp.add_apply,
        Nat.factorization_eq_zero_of_not_dvd
          (fun hd => hpn ((Nat.Prime.dvd_factorial hp).mp hd)),
        primorial_factorization_eq_zero hp hpn]
      exact ⟨0, rfl⟩

/-- **Witness layer (explicit square roots for the ⟸ direction of the main theorem)**: `x² = n! * n#`
    at each of n = 0..5, with x = 1, 1, 2, 6, 12, 60 (the pointwise correspondence of the main statement). -/
theorem witnesses_zero_to_five :
    1 ^ 2 = 0 ! * primorial 0 ∧
    1 ^ 2 = 1 ! * primorial 1 ∧
    2 ^ 2 = 2 ! * primorial 2 ∧
    6 ^ 2 = 3 ! * primorial 3 ∧
    12 ^ 2 = 4 ! * primorial 4 ∧
    60 ^ 2 = 5 ! * primorial 5 := by decide

/-- Bridge lemma (internal to the witness layer): an odd factorization at some prime implies non-square.
    `IsSquare` gives `m = r * r`, so `ν_p(m) = ν_p(r) + ν_p(r)` is even, contradicting the oddness assumption. -/
theorem not_isSquare_of_odd_factorization {m p e : ℕ} (hm : m ≠ 0)
    (h : m.factorization p = 2 * e + 1) : ¬ IsSquare m := by
  rintro ⟨r, hr⟩
  have hne : r ≠ 0 := by rintro rfl; rw [mul_zero] at hr; exact hm hr
  rw [hr, Nat.factorization_mul hne hne, Finsupp.add_apply] at h
  omega

/-- Non-square witness n = 6: `ν_2(6! * 6#) = 5` (odd); `6! * 6# = 2^5 * 675`, and `2 ∤ 675`. -/
theorem nonsquare_at_6 : ¬ IsSquare (6 ! * primorial 6) := by
  have hM : (6 : ℕ) ! * primorial 6 = 2 ^ 5 * 675 := by decide
  apply not_isSquare_of_odd_factorization (p := 2) (e := 2)
    (mul_ne_zero (Nat.factorial_ne_zero 6) (primorial_ne_zero 6))
  rw [hM, Nat.factorization_mul (by norm_num : ((2:ℕ)^5) ≠ 0) (by norm_num : (675:ℕ) ≠ 0),
    Finsupp.add_apply, Nat.factorization_eq_zero_of_not_dvd (by decide : ¬ (2:ℕ) ∣ 675),
    Nat.factorization_pow, Finsupp.smul_apply,
    Nat.Prime.factorization Nat.prime_two, Finsupp.single_eq_same]; norm_num [nsmul_eq_mul]

/-- Non-square witness n = 7: `ν_2(7! * 7#) = 5` (odd); `7! * 7# = 2^5 * 33075`, and `2 ∤ 33075`. -/
theorem nonsquare_at_7 : ¬ IsSquare (7 ! * primorial 7) := by
  have hM : (7 : ℕ) ! * primorial 7 = 2 ^ 5 * 33075 := by decide
  apply not_isSquare_of_odd_factorization (p := 2) (e := 2)
    (mul_ne_zero (Nat.factorial_ne_zero 7) (primorial_ne_zero 7))
  rw [hM, Nat.factorization_mul (by norm_num : ((2:ℕ)^5) ≠ 0) (by norm_num : (33075:ℕ) ≠ 0),
    Finsupp.add_apply, Nat.factorization_eq_zero_of_not_dvd (by decide : ¬ (2:ℕ) ∣ 33075),
    Nat.factorization_pow, Finsupp.smul_apply,
    Nat.Prime.factorization Nat.prime_two, Finsupp.single_eq_same]; norm_num [nsmul_eq_mul]

/-- Non-square witness n = 8: `ν_3(8! * 8#) = 3` (odd); `8! * 8# = 3^3 * 313600`, and `3 ∤ 313600`. -/
theorem nonsquare_at_8 : ¬ IsSquare (8 ! * primorial 8) := by
  have hM : (8 : ℕ) ! * primorial 8 = 3 ^ 3 * 313600 := by decide
  apply not_isSquare_of_odd_factorization (p := 3) (e := 1)
    (mul_ne_zero (Nat.factorial_ne_zero 8) (primorial_ne_zero 8))
  rw [hM, Nat.factorization_mul (by norm_num : ((3:ℕ)^3) ≠ 0) (by norm_num : (313600:ℕ) ≠ 0),
    Finsupp.add_apply, Nat.factorization_eq_zero_of_not_dvd (by decide : ¬ (3:ℕ) ∣ 313600),
    Nat.factorization_pow, Finsupp.smul_apply,
    Nat.Prime.factorization Nat.prime_three, Finsupp.single_eq_same]; norm_num [nsmul_eq_mul]

/-- Non-square witness n = 9: `ν_3(9! * 9#) = 5` (odd); `9! * 9# = 3^5 * 313600`, and `3 ∤ 313600`. -/
theorem nonsquare_at_9 : ¬ IsSquare (9 ! * primorial 9) := by
  have hM : (9 : ℕ) ! * primorial 9 = 3 ^ 5 * 313600 := by decide
  apply not_isSquare_of_odd_factorization (p := 3) (e := 2)
    (mul_ne_zero (Nat.factorial_ne_zero 9) (primorial_ne_zero 9))
  rw [hM, Nat.factorization_mul (by norm_num : ((3:ℕ)^5) ≠ 0) (by norm_num : (313600:ℕ) ≠ 0),
    Finsupp.add_apply, Nat.factorization_eq_zero_of_not_dvd (by decide : ¬ (3:ℕ) ∣ 313600),
    Nat.factorization_pow, Finsupp.smul_apply,
    Nat.Prime.factorization Nat.prime_three, Finsupp.single_eq_same]; norm_num [nsmul_eq_mul]

/-- Non-square witness n = 10: `ν_2(10! * 10#) = 9` (odd); `10! * 10# = 2^9 * 1488375`, and `2 ∤ 1488375`. -/
theorem nonsquare_at_10 : ¬ IsSquare (10 ! * primorial 10) := by
  have hM : (10 : ℕ) ! * primorial 10 = 2 ^ 9 * 1488375 := by decide
  apply not_isSquare_of_odd_factorization (p := 2) (e := 4)
    (mul_ne_zero (Nat.factorial_ne_zero 10) (primorial_ne_zero 10))
  rw [hM, Nat.factorization_mul (by norm_num : ((2:ℕ)^9) ≠ 0) (by norm_num : (1488375:ℕ) ≠ 0),
    Finsupp.add_apply, Nat.factorization_eq_zero_of_not_dvd (by decide : ¬ (2:ℕ) ∣ 1488375),
    Nat.factorization_pow, Finsupp.smul_apply,
    Nat.Prime.factorization Nat.prime_two, Finsupp.single_eq_same]; norm_num [nsmul_eq_mul]

/-- Non-square witness n = 11: `ν_2(11! * 11#) = 9` (odd); `11! * 11# = 2^9 * 180093375`, and `2 ∤ 180093375`. -/
theorem nonsquare_at_11 : ¬ IsSquare (11 ! * primorial 11) := by
  have hM : (11 : ℕ) ! * primorial 11 = 2 ^ 9 * 180093375 := by decide
  apply not_isSquare_of_odd_factorization (p := 2) (e := 4)
    (mul_ne_zero (Nat.factorial_ne_zero 11) (primorial_ne_zero 11))
  rw [hM, Nat.factorization_mul (by norm_num : ((2:ℕ)^9) ≠ 0) (by norm_num : (180093375:ℕ) ≠ 0),
    Finsupp.add_apply, Nat.factorization_eq_zero_of_not_dvd (by decide : ¬ (2:ℕ) ∣ 180093375),
    Nat.factorization_pow, Finsupp.smul_apply,
    Nat.Prime.factorization Nat.prime_two, Finsupp.single_eq_same]; norm_num [nsmul_eq_mul]

/-- Non-square witness n = 12: `ν_2(12! * 12#) = 11` (odd); `12! * 12# = 2^11 * 540280125`, and `2 ∤ 540280125`. -/
theorem nonsquare_at_12 : ¬ IsSquare (12 ! * primorial 12) := by
  have hM : (12 : ℕ) ! * primorial 12 = 2 ^ 11 * 540280125 := by decide
  apply not_isSquare_of_odd_factorization (p := 2) (e := 5)
    (mul_ne_zero (Nat.factorial_ne_zero 12) (primorial_ne_zero 12))
  rw [hM, Nat.factorization_mul (by norm_num : ((2:ℕ)^11) ≠ 0) (by norm_num : (540280125:ℕ) ≠ 0),
    Finsupp.add_apply, Nat.factorization_eq_zero_of_not_dvd (by decide : ¬ (2:ℕ) ∣ 540280125),
    Nat.factorization_pow, Finsupp.smul_apply,
    Nat.Prime.factorization Nat.prime_two, Finsupp.single_eq_same]; norm_num [nsmul_eq_mul]

/-- Non-square witness n = 13: `ν_2(13! * 13#) = 11` (odd); `13! * 13# = 2^11 * 91307341125`, and `2 ∤ 91307341125`. -/
theorem nonsquare_at_13 : ¬ IsSquare (13 ! * primorial 13) := by
  have hM : (13 : ℕ) ! * primorial 13 = 2 ^ 11 * 91307341125 := by decide
  apply not_isSquare_of_odd_factorization (p := 2) (e := 5)
    (mul_ne_zero (Nat.factorial_ne_zero 13) (primorial_ne_zero 13))
  rw [hM, Nat.factorization_mul (by norm_num : ((2:ℕ)^11) ≠ 0) (by norm_num : (91307341125:ℕ) ≠ 0),
    Finsupp.add_apply, Nat.factorization_eq_zero_of_not_dvd (by decide : ¬ (2:ℕ) ∣ 91307341125),
    Nat.factorization_pow, Finsupp.smul_apply,
    Nat.Prime.factorization Nat.prime_two, Finsupp.single_eq_same]; norm_num [nsmul_eq_mul]

/-- Non-square witness n = 14: `ν_5(14! * 14#) = 3` (odd); `14! * 14# = 5^3 * 20943712677888`, and `5 ∤ 20943712677888`. -/
theorem nonsquare_at_14 : ¬ IsSquare (14 ! * primorial 14) := by
  have hM : (14 : ℕ) ! * primorial 14 = 5 ^ 3 * 20943712677888 := by decide
  apply not_isSquare_of_odd_factorization (p := 5) (e := 1)
    (mul_ne_zero (Nat.factorial_ne_zero 14) (primorial_ne_zero 14))
  rw [hM, Nat.factorization_mul (by norm_num : ((5:ℕ)^3) ≠ 0) (by norm_num : (20943712677888:ℕ) ≠ 0),
    Finsupp.add_apply, Nat.factorization_eq_zero_of_not_dvd (by decide : ¬ (5:ℕ) ∣ 20943712677888),
    Nat.factorization_pow, Finsupp.smul_apply,
    Nat.Prime.factorization (by norm_num : Nat.Prime 5), Finsupp.single_eq_same]; norm_num [nsmul_eq_mul]

/-- Non-square witness n = 15: `ν_3(15! * 15#) = 7` (odd); `15! * 15# = 3^7 * 17955857920000`, and `3 ∤ 17955857920000`. -/
theorem nonsquare_at_15 : ¬ IsSquare (15 ! * primorial 15) := by
  have hM : (15 : ℕ) ! * primorial 15 = 3 ^ 7 * 17955857920000 := by decide
  apply not_isSquare_of_odd_factorization (p := 3) (e := 3)
    (mul_ne_zero (Nat.factorial_ne_zero 15) (primorial_ne_zero 15))
  rw [hM, Nat.factorization_mul (by norm_num : ((3:ℕ)^7) ≠ 0) (by norm_num : (17955857920000:ℕ) ≠ 0),
    Finsupp.add_apply, Nat.factorization_eq_zero_of_not_dvd (by decide : ¬ (3:ℕ) ∣ 17955857920000),
    Nat.factorization_pow, Finsupp.smul_apply,
    Nat.Prime.factorization Nat.prime_three, Finsupp.single_eq_same]; norm_num [nsmul_eq_mul]

/-- Non-square witness n = 16: `ν_3(16! * 16#) = 7` (odd); `16! * 16# = 3^7 * 287293726720000`, and `3 ∤ 287293726720000`. -/
theorem nonsquare_at_16 : ¬ IsSquare (16 ! * primorial 16) := by
  have hM : (16 : ℕ) ! * primorial 16 = 3 ^ 7 * 287293726720000 := by decide
  apply not_isSquare_of_odd_factorization (p := 3) (e := 3)
    (mul_ne_zero (Nat.factorial_ne_zero 16) (primorial_ne_zero 16))
  rw [hM, Nat.factorization_mul (by norm_num : ((3:ℕ)^7) ≠ 0) (by norm_num : (287293726720000:ℕ) ≠ 0),
    Finsupp.add_apply, Nat.factorization_eq_zero_of_not_dvd (by decide : ¬ (3:ℕ) ∣ 287293726720000),
    Nat.factorization_pow, Finsupp.smul_apply,
    Nat.Prime.factorization Nat.prime_three, Finsupp.single_eq_same]; norm_num [nsmul_eq_mul]

/-- Non-square witness n = 17: `ν_3(17! * 17#) = 7` (odd); `17! * 17# = 3^7 * 83027887022080000`, and `3 ∤ 83027887022080000`. -/
theorem nonsquare_at_17 : ¬ IsSquare (17 ! * primorial 17) := by
  have hM : (17 : ℕ) ! * primorial 17 = 3 ^ 7 * 83027887022080000 := by decide
  apply not_isSquare_of_odd_factorization (p := 3) (e := 3)
    (mul_ne_zero (Nat.factorial_ne_zero 17) (primorial_ne_zero 17))
  rw [hM, Nat.factorization_mul (by norm_num : ((3:ℕ)^7) ≠ 0) (by norm_num : (83027887022080000:ℕ) ≠ 0),
    Finsupp.add_apply, Nat.factorization_eq_zero_of_not_dvd (by decide : ¬ (3:ℕ) ∣ 83027887022080000),
    Nat.factorization_pow, Finsupp.smul_apply,
    Nat.Prime.factorization Nat.prime_three, Finsupp.single_eq_same]; norm_num [nsmul_eq_mul]

/-- Non-square witness n = 18: `ν_2(18! * 18#) = 17` (odd); `18! * 18# = 2^17 * 24936491397943125`, and `2 ∤ 24936491397943125`. -/
theorem nonsquare_at_18 : ¬ IsSquare (18 ! * primorial 18) := by
  have hM : (18 : ℕ) ! * primorial 18 = 2 ^ 17 * 24936491397943125 := by decide
  apply not_isSquare_of_odd_factorization (p := 2) (e := 8)
    (mul_ne_zero (Nat.factorial_ne_zero 18) (primorial_ne_zero 18))
  rw [hM, Nat.factorization_mul (by norm_num : ((2:ℕ)^17) ≠ 0) (by norm_num : (24936491397943125:ℕ) ≠ 0),
    Finsupp.add_apply, Nat.factorization_eq_zero_of_not_dvd (by decide : ¬ (2:ℕ) ∣ 24936491397943125),
    Nat.factorization_pow, Finsupp.smul_apply,
    Nat.Prime.factorization Nat.prime_two, Finsupp.single_eq_same]; norm_num [nsmul_eq_mul]

/-- Non-square witness n = 19: `ν_2(19! * 19#) = 17` (odd); `19! * 19# = 2^17 * 9002073394657468125`, and `2 ∤ 9002073394657468125`. -/
theorem nonsquare_at_19 : ¬ IsSquare (19 ! * primorial 19) := by
  have hM : (19 : ℕ) ! * primorial 19 = 2 ^ 17 * 9002073394657468125 := by decide
  apply not_isSquare_of_odd_factorization (p := 2) (e := 8)
    (mul_ne_zero (Nat.factorial_ne_zero 19) (primorial_ne_zero 19))
  rw [hM, Nat.factorization_mul (by norm_num : ((2:ℕ)^17) ≠ 0) (by norm_num : (9002073394657468125:ℕ) ≠ 0),
    Finsupp.add_apply, Nat.factorization_eq_zero_of_not_dvd (by decide : ¬ (2:ℕ) ∣ 9002073394657468125),
    Nat.factorization_pow, Finsupp.smul_apply,
    Nat.Prime.factorization Nat.prime_two, Finsupp.single_eq_same]; norm_num [nsmul_eq_mul]

/-- Non-square witness n = 20: `ν_2(20! * 20#) = 19` (odd); `20! * 20# = 2^19 * 45010366973287340625`, and `2 ∤ 45010366973287340625`. -/
theorem nonsquare_at_20 : ¬ IsSquare (20 ! * primorial 20) := by
  have hM : (20 : ℕ) ! * primorial 20 = 2 ^ 19 * 45010366973287340625 := by decide
  apply not_isSquare_of_odd_factorization (p := 2) (e := 9)
    (mul_ne_zero (Nat.factorial_ne_zero 20) (primorial_ne_zero 20))
  rw [hM, Nat.factorization_mul (by norm_num : ((2:ℕ)^19) ≠ 0) (by norm_num : (45010366973287340625:ℕ) ≠ 0),
    Finsupp.add_apply, Nat.factorization_eq_zero_of_not_dvd (by decide : ¬ (2:ℕ) ∣ 45010366973287340625),
    Nat.factorization_pow, Finsupp.smul_apply,
    Nat.Prime.factorization Nat.prime_two, Finsupp.single_eq_same]; norm_num [nsmul_eq_mul]

/-- Non-square witness n = 21: `ν_2(21! * 21#) = 19` (odd); `21! * 21# = 2^19 * 945217706439034153125`, and `2 ∤ 945217706439034153125`. -/
theorem nonsquare_at_21 : ¬ IsSquare (21 ! * primorial 21) := by
  have hM : (21 : ℕ) ! * primorial 21 = 2 ^ 19 * 945217706439034153125 := by decide
  apply not_isSquare_of_odd_factorization (p := 2) (e := 9)
    (mul_ne_zero (Nat.factorial_ne_zero 21) (primorial_ne_zero 21))
  rw [hM, Nat.factorization_mul (by norm_num : ((2:ℕ)^19) ≠ 0) (by norm_num : (945217706439034153125:ℕ) ≠ 0),
    Finsupp.add_apply, Nat.factorization_eq_zero_of_not_dvd (by decide : ¬ (2:ℕ) ∣ 945217706439034153125),
    Nat.factorization_pow, Finsupp.smul_apply,
    Nat.Prime.factorization Nat.prime_two, Finsupp.single_eq_same]; norm_num [nsmul_eq_mul]

/-- Assembled form (cited by the ⟹ direction of the main theorem): `n! * n#` is non-square at every `n` with `6 ≤ n ≤ 21`. -/
theorem nonsquare_six_to_twenty_one (n : ℕ) (h6 : 6 ≤ n) (h21 : n ≤ 21) :
    ¬ IsSquare (n ! * primorial n) := by
  interval_cases n
  · exact nonsquare_at_6
  · exact nonsquare_at_7
  · exact nonsquare_at_8
  · exact nonsquare_at_9
  · exact nonsquare_at_10
  · exact nonsquare_at_11
  · exact nonsquare_at_12
  · exact nonsquare_at_13
  · exact nonsquare_at_14
  · exact nonsquare_at_15
  · exact nonsquare_at_16
  · exact nonsquare_at_17
  · exact nonsquare_at_18
  · exact nonsquare_at_19
  · exact nonsquare_at_20
  · exact nonsquare_at_21

/-- **Lemma B (interval prime lemma)**: for every integer `n ≥ 22` there exists a prime `p` with `n/3 < p ≤ n/2`.
    (`/` is natural-number division: the interval reads `⌊n/3⌋ < p ≤ ⌊n/2⌋`; for prime existence this coincides
    with the real reading, since `p` is an integer: `n/3 < p ⟺ ⌊n/3⌋ < p` and `p ≤ n/2 ⟺ p ≤ ⌊n/2⌋`.) -/
theorem exists_prime_between_div_three_and_div_two (n : ℕ) (hn : 22 ≤ n) :
    ∃ p : ℕ, p.Prime ∧ n / 3 < p ∧ p ≤ n / 2 := by
  rcases lt_or_ge n 1000000 with h | h
  · exact NFactPrimorialSquareChain.lemmaB_upto_999999 n hn (by omega)
  · exact NFactPrimorialSquareChain.lemmaB_tail_1000000 n h

/-- **Main theorem (complete solution set, iff form)**:
    `n! * n#` is a perfect square ⟺ `n ≤ 5` (for `n : ℕ`).
    Over ℕ, `n ≤ 5 ⟺ n ∈ {0,1,2,3,4,5}`, so this is the main statement "the solutions are exactly n ∈ {0,...,5}". -/
theorem isSquare_factorial_mul_primorial_iff (n : ℕ) :
    IsSquare (n ! * primorial n) ↔ n ≤ 5 := by
  constructor
  · -- ⟹ direction: n! * n# a square ⟹ n ≤ 5 (contrapositive: for n ≥ 6, n! * n# is not a square)
    intro h
    by_contra hle
    have h6 : 6 ≤ n := by omega
    rcases lt_or_ge n 22 with hlt | h22
    · -- the 6 ≤ n ≤ 21 segment: direct contradiction with the pointwise non-square witnesses of the witness layer
      exact nonsquare_six_to_twenty_one n h6 (by omega) h
    · -- the 22 ≤ n segment: Lemma B yields a prime p ∈ (n/3, n/2]; the criterion gives Odd (ν_p(n!)), but Legendre's sum evaluates ν_p(n!) = 2
      obtain ⟨p, hp, hpl, hpu⟩ := exists_prime_between_div_three_and_div_two n h22
      have hodd : Odd ((n !).factorization p) :=
        (isSquare_iff_all_factorization_odd n (by omega)).mp h p hp (by omega)
      have hn3p : n < 3 * p := by omega
      have hp3 : 3 ≤ p := by omega
      have hpsq : n < p ^ 2 := by
        have h3pp : 3 * p ≤ p * p := Nat.mul_le_mul hp3 (le_refl p)
        rw [pow_two]
        exact lt_of_lt_of_le hn3p h3pp
      have hlog2 : Nat.log p n < 2 := Nat.log_lt_of_lt_pow (by omega) hpsq
      have hd1 : n / p = 2 := by
        have hp0 : 0 < p := hp.pos
        have hr : n - 2 * p < p := by omega
        rw [show n = n - 2 * p + p * 2 by omega, Nat.add_mul_div_left _ _ hp0,
          Nat.div_eq_of_lt hr, zero_add]
      have hfact2 : (n !).factorization p = 2 := by
        have hIco : Finset.Ico 1 2 = ({1} : Finset ℕ) := by decide
        calc (n !).factorization p = ∑ i ∈ Finset.Ico 1 2, n / p ^ i :=
              Nat.factorization_factorial hp hlog2
          _ = ∑ i ∈ ({1} : Finset ℕ), n / p ^ i := by rw [hIco]
          _ = n / p ^ 1 := Finset.sum_singleton _ _
          _ = n / p := by rw [pow_one]
          _ = 2 := hd1
      rw [hfact2] at hodd
      obtain ⟨k, hk⟩ := hodd
      omega
  · -- ⟸ direction: for n ≤ 5, give the explicit square root pointwise (x = 1, 1, 2, 6, 12, 60)
    intro h
    obtain ⟨h0, h1, h2, h3, h4, h5⟩ := witnesses_zero_to_five
    interval_cases n
    · exact (isSquare_iff_exists_sq _).mpr ⟨1, h0.symm⟩
    · exact (isSquare_iff_exists_sq _).mpr ⟨1, h1.symm⟩
    · exact (isSquare_iff_exists_sq _).mpr ⟨2, h2.symm⟩
    · exact (isSquare_iff_exists_sq _).mpr ⟨6, h3.symm⟩
    · exact (isSquare_iff_exists_sq _).mpr ⟨12, h4.symm⟩
    · exact (isSquare_iff_exists_sq _).mpr ⟨60, h5.symm⟩

/-- **Main theorem (existence-witness form, delivered alongside the iff form)**:
    the Diophantine equation `x² = n! * n#` has a solution ⟺ `n ∈ {0,1,2,3,4,5}`. -/
theorem exists_sq_factorial_mul_primorial_iff (n : ℕ) :
    (∃ x : ℕ, x ^ 2 = n ! * primorial n) ↔ n ∈ ({0, 1, 2, 3, 4, 5} : Finset ℕ) := by
  have hbridge : (∃ x : ℕ, x ^ 2 = n ! * primorial n) ↔ IsSquare (n ! * primorial n) := by
    rw [isSquare_iff_exists_sq]
    constructor
    · rintro ⟨x, hx⟩
      exact ⟨x, hx.symm⟩
    · rintro ⟨r, hr⟩
      exact ⟨r, hr.symm⟩
  rw [hbridge, isSquare_factorial_mul_primorial_iff]
  simp only [Finset.mem_insert, Finset.mem_singleton]
  constructor <;> intro h <;> omega

end NFactPrimorialSquare
\end{lstlisting}

\subsection{A refuted route, kept on the record}\label{sec:refuted}

The route originally envisaged for Lemma~\ref{lem:B} was a black-box splice
of Bertrand's postulate (mathlib's
\texttt{Nat.exists\_prime\_lt\_and\_le\_two\_mul}: a prime in $(m, 2m]$)
across two dyadic sub-intervals. During the task this route was not merely
found difficult but \emph{strictly refuted as a black-box argument}: the set
$\mathbb{P} \setminus \{13\}$ of all primes with $13$ removed still has the
Bertrand property --- every interval $(m, 2m]$ contains one of its members,
verified numerically for all $m \le 200000$ and closed by hand on
$m \in [7, 12]$ --- yet for $n = 33$ the target window $(11, 16]$ is empty in
that set, since with $13$ removed none of $12, 14, 15, 16$ is a member. The
window ratio $3/2$ is strictly below the Bertrand ratio $2$, so no
property-level argument at ratio $2$ can close it; the countermodel is a
statement about the proof method, not about the lemma (which is true, and is
proved here by the chain-and-tail route). We record this as an
honest negative result: the originally planned route was falsified, the
falsification is retained in the record, and the shipped proof does
not use it.

\section{Verification evidence}\label{sec:verify}

\begin{itemize}
\item \textbf{Toolchain.} Lean \texttt{v4.34.0}
(\texttt{leanprover/lean4:v4.34.0}), mathlib4 \texttt{v4.34.0}, revision
\texttt{5ed29652}, pinned in the artifact (\texttt{proofs/lean-toolchain}).
\item \textbf{Statement integrity.} The frozen snapshot
\texttt{nfact-primorial-square.statement.txt} has sha256
\begin{center}
\texttt{0b13d08d6d915c2ebb25a070461a9d2a}\allowbreak
\texttt{7d0364b12b3627c1b1efba5cc9f8ea48},
\end{center}
the statement source \texttt{nfact-primorial-square.lean} has sha256
\begin{center}
\texttt{638064201e96a3bdb42e75b00020c16e}\allowbreak
\texttt{2fb27c2738c58fa3d4923cde2a451a8f},
\end{center}
and the final proof file \texttt{final-nfact-primorial-square.lean}
(1094 lines, LF) has sha256
\begin{center}
\texttt{276cb2e698f6d201752eed3cb1104254}\allowbreak
\texttt{d976f1a392ba3860a194eaa66986070a}.
\end{center}
The final audit compared all eight statement headers of the proof file
against the frozen snapshot under two whitespace conventions (eight of eight
match) and re-checked the kernel goal states against the snapshot's
\texttt{\#check} record (eight of eight match). The listing excerpts in this
paper are byte-identical to the shipped proof files (independent mechanical
record: \texttt{audit/listing-verify.txt}). The three \texttt{proofs/} sources
are comment-cleaned editions (English docstrings): their code regions are
byte-identical to the frozen files listed above (correspondence:
\texttt{audit/comment-cleanup-diff.md}), and the post-cleanup kernel
re-verification is recorded in \texttt{audit/kernel-rerun.txt}.
\item \textbf{Well-definedness and non-degeneracy.} At the
well-definedness gate the statements were compiled with placeholder bodies
and probed with bare tactics: all nine probes fail, i.e.\ no single bare
tactic discharges the frozen statements (the gate output ships verbatim as
\texttt{audit/welldef-nfact-primorial-square.txt}). This is evidence about
the statements, not a proof of them.
\item \textbf{Kernel check.} A strict compile of the final file exits $0$
with no \texttt{sorry} and no \texttt{admit} (whole-file scan including
comments: zero hits; likewise zero hits for \texttt{axiom}, \texttt{unsafe},
\texttt{opaque}, \texttt{native\_decide}, \texttt{implemented\_by}). Three
diagnostics remain, all linter/deprecation-class and non-rejectable under
the audit discipline applied: an elaborator exponentiation-threshold hint on
\texttt{18\^{}100000} (intentional; the kernel checks the proof term
independently), a deprecated \texttt{push\_neg} alias, and one
unused-variable linter note. The axiom printout below is the verbatim output
of an independent re-run performed by the paper's packaging track on
2026-10-06 (strict compile exit $0$; the wrapped fourth entry is terminal
formatting; record: \texttt{audit/axioms.out}):
\begin{lstlisting}[basicstyle=\ttfamily\footnotesize,frame=none,label={lst:axioms}]
'NFactPrimorialSquare.isSquare_factorial_mul_primorial_iff' depends on axioms: [propext, Classical.choice, Quot.sound]
'NFactPrimorialSquare.exists_sq_factorial_mul_primorial_iff' depends on axioms: [propext, Classical.choice, Quot.sound]
'NFactPrimorialSquare.isSquare_iff_all_factorization_odd' depends on axioms: [propext, Classical.choice, Quot.sound]
'NFactPrimorialSquare.exists_prime_between_div_three_and_div_two' depends on axioms: [propext,
 Classical.choice,
 Quot.sound]
'NFactPrimorialSquare.witnesses_zero_to_five' depends on axioms: [propext, Classical.choice, Quot.sound]
'NFactPrimorialSquare.primorial_factorization_eq_one' depends on axioms: [propext, Classical.choice, Quot.sound]
'NFactPrimorialSquare.primorial_factorization_eq_zero' depends on axioms: [propext, Classical.choice, Quot.sound]
'NFactPrimorialSquare.factorization_factorial_mul_primorial' depends on axioms: [propext, Classical.choice, Quot.sound]
\end{lstlisting}
Every entry is exactly the whitelist \{propext, \texttt{Quot.sound},
\texttt{Classical.choice}\}; \texttt{sorryAx} does not occur.
\item \textbf{Grading.} The independent final audit's adjudication
(2026-10-06, four checks all passed) grades the eight frozen statements:
\emph{complete proof} --- zero \texttt{sorry}, axiom set inside the
whitelist, statement text unchanged from the frozen snapshot, strict compile
exit $0$. A cleaned copy of the adjudication ships as
\texttt{audit/gate3-final-adjudication.txt} (translated to English;
machine-local path prefixes replaced by \texttt{<REPO-ROOT>} and dated task
paths by \texttt{<source-task-id>}; verdict text and data unchanged ---
normalization disclosed in \texttt{audit/package-notes.md}).
\end{itemize}

\section{Related work and novelty evidence}\label{sec:related}

\paragraph{The source problem and its neighbours.} The equation is one of
two cases recorded as open in Novakovi\'c's paper \cite{novakovic2026}
(quotation in Section~\ref{sec:intro}); that paper studies
$P(x) = \prod_{i=1}^{r} H_{n_i}$, a generalization of the Brocard--Ramanujan
equation. The Brocard-side computational bound ($10^{9}$, attributed in the
source task's records to Berndt and Galway) is background only: the exclusion
of $n \ge 22$ in this paper is a pure theorem (Lemma~\ref{lem:B}) and does
not rely on it; that reference is recorded as undecided at the full-text
level in the source task and is not load-bearing here. The same author's
companion paper \cite{novakovic2026b} is recorded with the restricted
wording that its object of study does not contain $n\#$.

\paragraph{Delimitations (wording obligations of the source task, discharged
here one by one).}
(i) OEIS A055204 \cite{oeisA055204} asks for the \emph{minimal} $a$ making
$n! \cdot a$ a square; the present theorem asks the \emph{specific} value
$a = n\#$ and is not an instance of A055204.
(ii) Erd\H{o}s problem \#399 \cite{erdos399} concerns $n! = x^{k} \pm y^{k}$
--- a different equation from $x^2 = n! \cdot n\#$ --- and is marked
DISPROVED (2026-02-04, counterexample $10! = 48^{4} - 36^{4}$); it is
recorded here only as same-family adjacency, never as an open anchor.
(iii) Peixoto \cite{peixoto2026} (\S 5.3.1) treats the primorial-core Pell
equation $X^{2} - D Z^{2} = 1$ --- an adjacent object that does not
characterize which $n$ make $n! \cdot n\#$ a square.
(iv) Yudin \cite{yudin2026} (square products of factorials and a conjecture
of Erd\H{o}s--Graham) is a different object; no collision, and the citation
duty is discharged here.
(v) Nagura \cite{nagura1952} (for $x \ge 25$ there is a prime in
$(x, 6x/5]$) and OEIS A166968 \cite{oeisA166968} (whose $a(2) = 8$ says: for
$m \ge 8$ there is a prime in $(m, 3m/2]$) provide a literature-level
shortest path toward Lemma~\ref{lem:B} in the $m$-parametrization; the
kernel proof shipped here does not go through it. Nagura's original text was
not obtained; its statement is corroborated by several OEIS
entries, and the load-bearing proof here is independent of it.
(vi) The failure set of the interval property below $n = 22$ is
$\{2,3,9,21\}$ under the domain convention $n \ge 2$ ($n = 0, 1$ also fail,
on an empty interval; this is irrelevant to Lemma~\ref{lem:B}, which starts
at $n = 22$); the failure set was confirmed by three independent numerical
recomputations.

\paragraph{Novelty evidence.} An exclusivity review adjudicated:
\textbf{no occupying record found};
search-strength label \textbf{strong}, with closing dates 2026-10-05
(deep-dig layer) and 2026-10-06 (an independent review layer that re-compiled
the load-bearing kernels, recomputed the numeric evidence from scratch and
maintained every load-bearing claim). A zero hit is not novelty; the
strongest statement licensed here is that no record of this object was found
in the channels queried. The full query log, including every zero-hit query,
ships as \texttt{audit/c9-record.md} in the deposit (English cleaned edition;
conclusions and query data unedited). No ``first formalization'' statement is
made anywhere in this paper.

\begin{center}\small
\begin{tabular}{p{3.4cm}p{5.4cm}p{6.0cm}}
\hline
channel & queries & result \\
\hline
literature (author self-report) & arXiv:2601.16757 body text & the case is on
record as open (quotation in Section~\ref{sec:intro}) \cite{novakovic2026} \\
arXiv metadata & primorial $\wedge$ factorial $\wedge$ perfect square;
primorial $\wedge$ factorial $\wedge$ square; double factorial $\wedge$
least common multiple; squarefree part $\wedge$ factorial & zero hits each \\
arXiv metadata & Brocard $\wedge$ primorial & one hit \cite{peixoto2026}:
adjacent object (primorial-core Pell), does not characterize this equation \\
OpenAlex full text & brocard/primorial/square/factorial family; k-free
Brocard family & nearest hits include
\cite{novakovic2026,novakovic2026b}; no same-form statement \\
OEIS keyword & ``n! primorial square''; ``factorial times primorial is a
square'' & hits exist; no same-form \\
OEIS keyword & ``double factorial lcm square''; ``Novakovic factorial'' &
zero hits \\
public Lean ecosystem (formal-conjectures tarball, content level) &
\texttt{primorial}; \texttt{doubleFactorial} & occurrences in unrelated
files only; no item in this domain \\
OEIS value columns (review layer) & $x$: 1,1,2,6,12,60; $x^2$: 1,4,36,144,3600
& no registration of these objects \\
local mathlib (pinned v4.34.0), positive control & content scan with live
control & channel live; no same-form item \\
\hline
\end{tabular}
\end{center}

\section{Limitations and threats to validity}\label{sec:limits}

(a) \textbf{What is proved.} The kernel layer proves exactly the eight frozen
statements of Listing~\ref{lst:stmt}; the grade of each is \emph{complete
proof}. Nothing beyond is claimed: no growth statement, no other multiplier
in place of $n\#$, and no second equation.

(b) \textbf{The second case is not claimed.} The companion equation
$x^2 = n!! \cdot [1, \ldots, n]$ of the same remark in \cite{novakovic2026}
is not claimed: only an exclusion path for even $n$ was obtained, and it is
recorded as a follow-up item.

(c) \textbf{No registration claim.} No OEIS registration is claimed or
performed; the review layer found no registration of the value columns, and
any registration action is the author's own.

(d) \textbf{Independence from computational bounds.} The exclusion of
$n \ge 22$ is a pure theorem and does not rely on the Brocard-side
computational bound ($10^{9}$), which is full-text-undecided background and
not load-bearing.

(e) \textbf{A stronger truth than the statement.} The hypothesis $n \ge 1$ of
Lemma~\ref{lem:A} is not used by its proof (the criterion holds at $n = 0$ as
well); the lemma is stated as frozen, and we report the fact rather than
silently strengthening the statement.

(f) \textbf{Comment-level deltas.} Three docstrings in the final file each
carry one added explanatory sentence relative to the frozen snapshot (a
notation note, the natural-division reading, and the correspondence to the
informal statement); docstrings are comments and do not enter the statement
comparison, which matched under both whitespace conventions.

(g) \textbf{Semantic fidelity of the formalization.} The theorems are about
the Lean definitions. Their fidelity to the informal reading rests on the
frozen statements being the intended encoding; a semantic round-trip check
(formal statement back to natural language, adjudicated consistent) was
performed at the formalization stage, and a further spot check of this
paper's informal exposition against the frozen statements was done at
assembly. A reader who disagrees with the encoding is disagreeing with the
statement, not with the proof.

(h) \textbf{Nagura's original text.} Not obtained; the statement is used only
as corroborated background (Section~\ref{sec:related}(v)) and is not
load-bearing.

(i) \textbf{Novelty review limits.} The zero-hit verdict is a statement about
the channels reached (arXiv metadata, OpenAlex full text, OEIS keyword and
value faces, the public Lean ecosystem via a repository tarball, the local
mathlib tree), not about the literature as a whole; a list of background
items recorded as undecided at the full-text level is enumerated in the
source task card (which stays in the source repository and is not shipped),
and none of them is load-bearing for this result. No ``first formalization''
claim is made.

(j) \textbf{Numeric evidence is layered.} The failure set, the chain
coverage and the A166968 correspondence are numeric- or literature-level
evidence (the recomputation counts are stated where used); every
theorem-level conclusion is adjudicated by the kernel and depends on no
numeric evidence.

\section*{Data availability}

The artifact deposited with this paper contains: the paper source and its PDF
(\texttt{paper/}); the Lean sources (\texttt{proofs/}: the frozen statement
snapshot, the statement source, the final proof file and the toolchain pin);
the audit evidence (\texttt{audit/}: the final adjudication, the
well-definedness gate output, the axiom printout, the listing verification
record, the bibliography verification record, the full novelty-search record, the
semantic fidelity record, and the package self-description); and the
metadata. Reproduction: with the pinned toolchain and mathlib revision, a
strict compile of the deposited proof file exits $0$, and appending the eight
\texttt{\#print axioms} lines reproduces the printout of
Section~\ref{sec:verify}. The companion claim of record for the same kernel
result is deposited separately \cite{claimdeposit} (concept DOI
\texttt{10.5281/zenodo.23177003}). This paper's own deposit is published on
Zenodo (DOI \texttt{10.5281/zenodo.23179326}); the version history is
recorded in the deposit's metadata README. The upload steps themselves are
the author's (\texttt{UPLOAD.md} in the source repository lists them; it is
not part of the deposit). Paper under CC BY 4.0, code under MIT.

\section*{Use of generative AI}

The formalization and proofs reported in this paper were developed with AI
assistance (large language model: \texttt{kimi-k3}) under human direction
and review. All statements and proofs reported here were checked by the Lean
4 kernel (Lean v4.34.0, mathlib v4.34.0): every proof compiles with no
\texttt{sorry}, and \texttt{\#print axioms} reports only subsets of
$\{$\texttt{propext}, \texttt{Quot.sound}, \texttt{Classical.choice}$\}$.
The author reviewed all content and takes full responsibility for it.

\section*{Author contributions}
The author selected and adjudicated the topic, verified the semantics of the
formal statements against intent, signed off the verification gates, and made
the submission decision. The AI system is not an author.

\begin{thebibliography}{9}
% Only works actually cited; each was re-checked to exist before entry
% (2026-10-06: lean4/mathlib/nagura1952 via live Crossref title (and pages)
% metadata; the four arXiv items via the live abs-page citation_title/
% citation_date metadata; erdos399 via the live page, whose status widget
% reads DISPROVED; the two OEIS entries via the live OEIS search name fields;
% claimdeposit via a live doi.org resolution). No entry is quoted from memory.
% Record: audit/bib-verify.md (a copy ships in the deposit).
\bibitem{lean4}
L.~de Moura and S.~Ullrich.
\newblock The Lean 4 theorem prover and programming language.
\newblock \emph{Automated Deduction -- CADE-28}, LNCS, pp.~625--635, 2021.
\url{https://doi.org/10.1007/978-3-030-79876-5_37}

\bibitem{mathlib}
The mathlib Community.
\newblock The Lean mathematical library.
\newblock \emph{CPP 2020}, pp.~367--381, 2020.
\url{https://doi.org/10.1145/3372885.3373824}

\bibitem{novakovic2026}
S.~Novakovi\'c.
\newblock The Diophantine equation $P(x) = \prod_{i=1}^{r} H_{n_i}$.
\newblock arXiv:2601.16757, 2026. \url{https://arxiv.org/abs/2601.16757}

\bibitem{novakovic2026b}
S.~Novakovi\'c.
\newblock On a generalization of the Brocard--Ramanujan Diophantine equation.
\newblock arXiv:2602.09687, 2026. \url{https://arxiv.org/abs/2602.09687}

\bibitem{peixoto2026}
L.~V. Peixoto.
\newblock The Brocard problem: structural invariants, p-adic density, and the
generative sieve.
\newblock arXiv:2606.23485, 2026. \url{https://arxiv.org/abs/2606.23485}

\bibitem{yudin2026}
F.~Yudin.
\newblock Square products of factorials and a conjecture of Erd\H{o}s and
Graham.
\newblock arXiv:2610.01899, 2026. \url{https://arxiv.org/abs/2610.01899}

\bibitem{nagura1952}
J.~Nagura.
\newblock On the interval containing at least one prime number.
\newblock \emph{Proceedings of the Japan Academy} 28 (1952), 177--181.
\url{https://doi.org/10.3792/pja/1195570997}

\bibitem{erdos399}
Erd\H{o}s Problems website.
\newblock Problem \#399 ($n! = x^k \pm y^k$; marked DISPROVED, counterexample
$10! = 48^4 - 36^4$).
\newblock \url{https://www.erdosproblems.com/399}

\bibitem{oeisA055204}
OEIS Foundation.
\newblock Squarefree part of $n!$: $n!$ divided by its largest square divisor
(sequence A055204).
\newblock \emph{The On-Line Encyclopedia of Integer Sequences}.
\url{https://oeis.org/A055204}

\bibitem{oeisA166968}
OEIS Foundation.
\newblock Minimum $k$ such that for all $m \ge k$ there is a prime $p$ with
$m < p < m(n+1)/n$ (sequence A166968).
\newblock \emph{The On-Line Encyclopedia of Integer Sequences}.
\url{https://oeis.org/A166968}

\bibitem{claimdeposit}
Aurora0134.
\newblock Priority claim: complete solution set of $x^2 = n! \cdot n\#$
(exactly $n \in \{0,\ldots,5\}$) --- Lean 4 formalized proof and claim note,
v1.0.2.
\newblock \emph{Zenodo}, 2026. Concept DOI:
\url{https://doi.org/10.5281/zenodo.23177003}
\end{thebibliography}

\end{document}
